\honest{Conservation of Expected Evidence}{Eliezer Yudkowsky}{2007}

In 1631 the priest Friedrich Spee described how accused witches were judged. A bad life proved guilt, and so did a good one, because witches pretend to be virtuous. Fear in prison proved guilt, and so did calm. Spee heard many confessions, so he could see every branch: whatever the woman did, it counted against her. \nb{Spee saw the point without any probability theory.}

It is for this reason, I say, that scientists write down their predictions in advance. \nb{In 2007, outside clinical trials, few scientists registered their predictions publicly before collecting data.}

Then my law, which I name ``Conservation of Expected Evidence'': your expected belief after seeing the evidence must equal your belief before. I write out three lines. \nb{The third is the law of total probability, taught in any first course in the subject.} Then comes the slogan: ``for every expectation of evidence, there is an equal and opposite expectation of counterevidence.''

A confident theory, I say, gains little when its prediction comes true, and loses a lot when the prediction fails. I end the paragraph: ``On average, you must expect to be exactly as confident as when you started out.'' \nb{Here ``confident'' means the probability of the theory. You do expect to end up more sure, one way or the other.}

Three applications follow. If no sabotage is evidence for a Fifth Column, then sabotage must be evidence against one. If a good life shows that a woman is a witch, then an evil life must show that she is not. Then: if God refuses to reveal Himself in order to test our faith, ``then the miracles described in the Bible must argue against the existence of God.'' ``Doesn't quite sound right, does it?'' I ask, and tell you to pay attention to that feeling. \nb{The last one follows only if God's silence is offered as evidence for God. The argument as stated only explains the silence, which would make miracles weaker evidence for God, not evidence against.}

Then I say: ``For a true Bayesian, it is impossible to seek evidence that confirms a theory.'' \nb{In plain English you can: a few paragraphs earlier the post described a test a theory will probably pass. The next sentence narrows the claim to raising your confidence on average.}

This, I say, ``can take quite a load off your mind,'' since you need not plan how to make the evidence confirm your theory. I advise: ``You might as well sit back and relax while you wait for the evidence to come in.'' \nb{That is advice for an ideal reasoner. A human can reinterpret results to suit a theory, and often does.}

I add: ``Human psychology is so screwed up.''
